% Part: sets-functions-relations
% Chapter: sets
% Section: subsets

\documentclass[../../../include/open-logic-section]{subfiles}

\begin{document}

\olfileid{sfr}{set}{sub}
\olsection{فرعي سټونه او د فرعي سټونو سټونه}

\begin{explain}
مونږ به ډېر ځله سټونه سره پرتله کول غواړو۔ د پرتله کولو يو څرګند ډول داسې دے: \emph{هر هغه څيز چې په يو سټ کښې وي، په بل کښې هم وي}۔ دا حالت دومره مهم دے چې د دې دپاره نوې نښې معرفي کړو۔
\end{explain}

\begin{defn}[فرعي سټ]
که د يو سټ $A$ هر !!{element} د~$B$ هم !!a{element} وي، نو وايو چې $A$ د~$B$ يو \emph{فرعي سټ} دے، او $A \subseteq B$ ليکو۔ که $A$ د~$B$ فرعي سټ نۀ وي، $A \not\subseteq B$ ليکو۔ که $A \subseteq B$ وي خو $A \neq B$ وي، $A \subsetneq B$ ليکو او وايو چې $A$ د $B$ يو \emph{خاص فرعي سټ} دے۔
\end{defn}

\begin{ex}
هر سټ د خپل ځان فرعي سټ دے، او $\emptyset$ د هر سټ فرعي سټ دے۔ د جفتو طبيعي عددونو سټ د طبيعي عددونو د سټ فرعي سټ دے۔ همدارنګه، $\{ a, b \} \subseteq \{ a, b, c \}$۔ خو $\{ a, b, e \}$ د $\{ a, b, c \}$ فرعي سټ نۀ دے۔
\end{ex}

\begin{ex}
عدد $2$ د صحيح عددونو د سټ يو !!{element} دے، حال دا چې د جفتو عددونو سټ د صحيح عددونو د سټ فرعي سټ دے۔ خو يو سټ کېدے شي د بل سټ \emph{هم} !!a{element} وي او هم فرعي سټ، د بېلګې په توګه، $\{0\} \in \{0, \{0\}\}$ او هم $\{0\} \subseteq \{0, \{0\}\}$۔
\end{ex}

د غړو له مخې برابري د سټونو د برابرۍ معيار راکوي: $A = B$ هله او يوازې هله وي چې د~$A$ هر !!{element} د~$B$ هم !!a{element} وي، او برعکس۔ د ``فرعي سټ'' تعريف $A \subseteq B$ بالکل د دې معيار د لومړۍ نيمې په توګه تعريفوي: د~$A$ هر !!{element} د~$B$ هم !!a{element} وي۔ البته، که $A$ او $B$ سره بدل کړو، تعريف بيا هم عملي کېږي: يعنې $B \subseteq A$ هله او يوازې هله وي چې د~$B$ هر !!{element} د~$A$ هم !!a{element} وي۔ او دا بيا بالکل د غړو له مخې د برابرۍ ``برعکس'' برخه ده۔ په نورو ټکو، د غړو له مخې له برابرۍ نه دا نتيجه راوځي چې سټونه هله او يوازې هله برابر وي چې د يو بل فرعي سټونه وي۔

\begin{prop}
$A = B$ هله او يوازې هله وي چې $A \subseteq B$ او $B \subseteq A$ دواړه وي۔
\end{prop}

اوس د ځينو نورو ګټورو نښو د معرفي کولو دپاره هم ښه موقع ده۔ د دې په تعريف کښې چې $A$ کله د~$B$ فرعي سټ وي، مونږ ووئيل چې ``د~$A$ هر !!{element} \dots دے''، او د ``$\dots$'' ځاے مو په ``د $B$ !!a{element}'' ډک کړ۔ خو دا د عبارت دومره عامه \emph{بڼه} ده چې د دې دپاره د صوري نښو معرفي کول به ګټور وي۔

\begin{defn}\ollabel{forallxina}
$(\forall x \in A)\phi$ د $\forall x(x \in A \lif \phi)$ لنډه بڼه ده۔ همدارنګه، $(\exists x \in A)\phi$ د $\exists x(x \in A \land \phi)$ لنډه بڼه ده۔
\end{defn}

د دې نښو په کارولو سره وئيلے شو چې $A \subseteq B$ هله او يوازې هله وي چې $(\forall x \in A)x \in B$ وي۔

اوس د يو ځانګړي ډول سټ څېړنې ته ځو: د يو ټاکلي سټ د ټولو فرعي سټونو سټ۔

\begin{defn}[د فرعي سټونو سټ]
هغه سټ چې د يو سټ~$A$ ټول فرعي سټونه پکښې وي، د~$A$ \emph{د فرعي سټونو سټ} بللے شي، او $\Pow{A}$ ليکل کېږي۔
  \[
    \Pow{A} = \Setabs{B}{B \subseteq A}
  \]
\end{defn}

\begin{ex}
د $\{ a, b, c \}$ ټول ممکن فرعي سټونه کوم دي؟ دا دي: $\emptyset$، $\{a \}$، $\{b\}$، $\{c\}$، $\{a, b\}$، $\{a, c\}$، $\{b, c\}$، $\{a, b, c\}$۔ د دې ټولو فرعي سټونو سټ $\Pow{\{a,b,c\}}$ دے:
\[
\Pow{\{ a, b, c \}} = \{\emptyset, \{a \}, \{b\}, \{c\}, \{a, b\},
\{b, c\}, \{a, c\}, \{a, b, c\}\}
\]
\end{ex}

\begin{prob}
د $\{a, b, c, d\}$ ټول فرعي سټونه وليکئ۔
\end{prob}

\begin{prob}
وښايئ چې که $A$ د $n$ په شمېر !!{element}s ولري، نو $\Pow{A}$ د $2^n$ په شمېر !!{element}s لري۔
\end{prob}

\end{document}
